\section{Sources and how they were checked}\label{app:sources}

The brief for this document restricted sources to arXiv, the IACR ePrint archive, and the published material of established quantum-computing companies, with the additional instruction to consult IBM Quantum's resources. The policy applied:

\begin{itemize}[leftmargin=*]
\item \textbf{arXiv and ePrint.} Every paper cited with an arXiv or ePrint identifier was located through its arXiv abstract page (\code{arxiv.org/abs/\ldots}) or its ePrint record, and the numbers quoted from it were taken from the abstract or from the paper's own tables (fetched as the arXiv HTML rendering where the abstract was not enough: \cite{babbush26,haener26,yoder25tour,cain26}). Two papers were read from the user's own copies in \code{shor/papers/}: Schrottenloher \cite{schrottenloher26} (ePrint 2026/1128) and Kim et al.\ \cite{kim26} (ePrint 2026/106).
\item \textbf{Company material.} A source counts as a company source when the paper's arXiv record lists the company as the authors' affiliation or the page is served from the company's own domain. The companies whose material is used, and the form it takes:
  \begin{description}[leftmargin=*,style=nextline]
  \item[Google Quantum AI (\code{quantumai.google}).] Willow below-threshold memory \cite{willow25}; colour code on a superconducting processor \cite{lacroix24color}; AlphaQubit (with Google DeepMind) \cite{alphaqubit24}; yoked surface codes \cite{yoked23}; magic state cultivation \cite{cultivation24}; the colour-code circuits and Chromobius \cite{gidney23color}; the RSA-2048 estimate \cite{gidney25rsa}; the ECDLP-256 estimates and mitigations paper \cite{babbush26}; Stim and PyMatching \cite{gidney21stim,higgott23pymatching}.
  \item[IBM Quantum (\code{ibm.com/quantum}).] The bivariate bicycle (gross) code paper \cite{bravyi24gross}; the bicycle architecture (``Tour de gross'') \cite{yoder25tour}; improved QLDPC surgery \cite{cross24surgery}; Relay-BP \cite{mueller25relaybp}; the June 2025 fault-tolerance announcement \cite{ibm-blog25} and the development roadmap \cite{ibm-roadmap}, both read from \code{ibm.com}; and the IBM Quantum Learning course ``Foundations of quantum error correction'' \cite{ibm-learning-qec}, used as the reference for the pedagogical ordering of \cref{sec:qec-why,sec:codes}.
  \item[IonQ (\code{ionq.com}).] The walking-cat architecture \cite{walkingcat26} and the ECDLP-256 trapped-ion estimate \cite{haener26} (both arXiv, IonQ affiliations; the latter is also ePrint 2026/1916).
  \item[Quantinuum / Honeywell (\code{quantinuum.com}).] Real-time fault-tolerant error correction on the $\bits{7}{1}{3}$ code \cite{ryan-anderson21} and fault-tolerant logical teleportation with transversal gates and lattice surgery \cite{ryan-anderson24}.
  \item[PsiQuantum (\code{psiquantum.com}).] The active-volume architecture \cite{litinski22active} and the ``50 million Toffoli gates'' ECDLP estimate \cite{litinski23}.
  \item[QuEra / Harvard (\code{quera.com}).] The logical processor on reconfigurable atom arrays \cite{bluvstein24} and logical magic-state distillation \cite{rodriguez24distill}; the 2026 neutral-atom Shor estimate \cite{cain26} is from the same Harvard/Caltech group with QuEra co-authors.
  \item[Alice\,\&\,Bob (\code{alice-bob.com}).] The repetition-cat ECDLP-256 estimate \cite{gouzien23} and the repetition-cat proposal \cite{guillaud19}.
  \end{description}
  Iceberg Quantum's Pinnacle architecture \cite{pinnacle26} is included in \cref{tab:external-phys} as an arXiv paper, not as an established-company source.
\item \textbf{The \ecdsafail{} material} is on GitHub: the challenge repository, the supplemental materials of the paper's authors, and the per-commit wiki (\cref{sec:challenge}); the paper itself is on arXiv \cite{ecdsafail26}.
\item \textbf{Textbook material} (stabiliser formalism, the small codes, the threshold theorem, the bounds) is cited to the original arXiv/journal papers \cite{gottesman97,crss98,knill-laflamme97,aharonov-benor08,shor95,steane96,laflamme96,bennett96,kitaev03,dennis02} and to Gottesman's textbook draft \cite{gottesman24book}.
\item \textbf{Computed for this document}: the encoders and syndrome tables of \cref{sec:smallcodes}, the surface-code threshold data of \cref{fig:threshold}, the BB-code verification of \cref{tab:bb}, the BP+OSD simulation of \cref{fig:bposd}, the bounds of \cref{fig:bounds}, the Pareto and trajectory data of \cref{sec:pareto}, and the Model S table. Every such number is marked ``computed'' at the point of use and reproducible from \cref{app:scripts}.
\end{itemize}
